\chapter{Algebraic Algorithms}\label{ch:algebraic}

\begin{center}\itshape
Whatever can be done with numbers can be done with polynomials, faster;\\
and multiplication is the engine that drives everything else.\\
\upshape(the creed of computer algebra, in our words)
\end{center}

In 1971 Arnold Sch\"onhage and Volker Strassen published an algorithm that multiplies two $n$-bit integers with $O(n\log n\log\log n)$ bit operations, and they conjectured that the
best possible bound should be $O(n\log n)$. For almost half a century the conjecture stood, and the $\log\log n$ was whittled down by ever more ingenious methods without disappearing.
On 18 March 2019 David Harvey, of the University of New South Wales, and Joris van der Hoeven, of the CNRS in France, posted a paper with the title \emph{Integer multiplication in time
$O(n\log n)$}. The algorithm is of purely theoretical interest, as its advantage only appears for numbers far larger than anything a computer could store; but the question that Kolmogorov
had opened in his Moscow seminar of 1960 (Chapter~\ref{ch:measuring}) now has, most experts believe, its final answer.

This chapter is about what happens once multiplication is fast. Tome~4, Chapter~7, proved the fast Fourier transform and its exact version modulo a prime, and Tome~S1, Chapter~10, proved that
Strassen's seven products give matrix multiplication in $O(n^{2.81})$. We start below the FFT, with Karatsuba's and Toom's splitting methods, then extend the FFT to rings where the roots of
unity are powers of $2$, which is the idea of Sch\"onhage and Strassen. Then multiplication is put to work. Newton's method, run on power series rather than on real numbers, divides as fast as it
multiplies. A tree of products evaluates a polynomial at $n$ points, and interpolates it back, in nearly linear time. Matrix inversion costs no more than matrix multiplication. And the $n$th term of a
linear recurrence comes from a polynomial computed by repeated squaring.

Throughout, $\mathsf M(n)$ denotes the number of ring operations needed to multiply two polynomials of degree less than $n$ over the ring at hand. By Tome~4, Corollary~7.10 and Theorem~7.11, $\mathsf M(n)=O(n\log n)$
over $\mathbb C$ and over the fields $\mathbb F_p$ with $p\equiv1\pmod{2^r}$ for $2^r\ge2n$; and $\mathsf M(n)=O(n\log n\log\log n)$ over any commutative ring, by Cantor and Kaltofen (1991), a result we only quote. We
assume that $\mathsf M(n)/n$ is nondecreasing, which all these bounds satisfy; then $\mathsf M(n)+\mathsf M(n/2)+\mathsf M(n/4)+\dots\le2\mathsf M(n)$.

\begin{think}
Dividing $1$ by $1-x$ as power series gives $1+x+x^2+\cdots$. If you know the first $k$ coefficients of $1/f$ for some power series $f$, how many coefficients can one multiplication give you? (Newton's method doubles the number of correct digits of a square root at each step; what plays the role of digits here?)
\end{think}

\section{Splitting: Karatsuba and Toom}\label{sec:toom}

Karatsuba's idea (Chapter~\ref{ch:measuring}, Section~\ref{sec:master}) multiplies $a_0+a_1X$ by $b_0+b_1X$ with three products, $a_0b_0$, $a_1b_1$ and $(a_0+a_1)(b_0+b_1)$, from which the middle coefficient is $(a_0+a_1)(b_0+b_1)-a_0b_0-a_1b_1$. Applied recursively to polynomials of degree less than
$n$, split into halves $A=A_0+X^{n/2}A_1$, it needs $\mathsf T(n)=3\mathsf T(n/2)+O(n)$ operations, so $O(n^{\log_23})$ by the master theorem (Theorem~\ref{thm:master}); Tome~4, Exercise~7.3, counts the multiplications exactly. For integers, the halves are the high and low digits, and the
only new point is that $a_0+a_1$ may have one more digit (Exercise~\ref{ex:karatsubaint}).

Andrei Toom saw in 1963 that Karatsuba's trick is evaluation and interpolation in disguise: $a_0b_0$, $a_1b_1$ and $(a_0+a_1)(b_0+b_1)$ are the values of the product at $0$, at ``infinity'' (the leading coefficient) and at $1$. Split each factor into $k$ pieces instead of two.

\begin{theorem}[Toom, 1963; Cook, 1966]\label{thm:toomcook}
Let $R$ be a commutative ring containing $2k-2$ distinct elements $x_1,\dots,x_{2k-2}$ whose pairwise differences are invertible. Polynomials of degree less than $n$ over $R$ can be multiplied with $O(n^{\log_k(2k-1)})$ operations of $R$, the constant depending on $k$.
\end{theorem}
\begin{proof}
Let $n=k^r$ (pad with zeros otherwise) and write $A=\sum_{i<k}A_iX^{in/k}$ and $B=\sum_{i<k}B_iX^{in/k}$ with $\deg A_i,\deg B_i<n/k$. Consider $a(Y)=\sum_iA_iY^i$ and $b(Y)=\sum_iB_iY^i$, polynomials in a new variable $Y$ of degree less than $k$ whose coefficients are polynomials in $X$. Their product
$c(Y)=\sum_{i<2k-1}C_iY^i$ has $2k-1$ unknown coefficients, and $AB=c(X^{n/k})$, which is assembled from the $C_i$ by $O(n)$ additions. To find the $C_i$, compute the $2k-1$ products $a(x_j)b(x_j)=c(x_j)$ for $j\le2k-2$, and $A_{k-1}B_{k-1}=C_{2k-2}$, the value ``at infinity''. Each $a(x_j)$ is a fixed linear combination of the
$A_i$, costing $O(n)$ operations for fixed $k$, and each product is a product of polynomials of degree less than $n/k$, done recursively. Knowing $C_{2k-2}$, the polynomial $c(Y)-C_{2k-2}Y^{2k-2}$ has degree at most $2k-3$ and known values at the $2k-2$ points $x_j$, so its coefficients are obtained by Lagrange interpolation, a fixed linear map whose coefficients are
products of the $(x_i-x_j)^{-1}$ (Tome~0, Chapter~3), applied coefficientwise in $O(n)$ operations. So $\mathsf T(n)=(2k-1)\mathsf T(n/k)+O(n)$, and since $\log_k(2k-1)>1$, the master theorem gives $O(n^{\log_k(2k-1)})$.
\end{proof}

Toom-$3$, with the points $0,1,-1,2$ and infinity, gives $n^{\log_35}\approx n^{1.465}$ (Exercise~\ref{ex:toom3}). As $k$ grows, the exponent $\log_k(2k-1)$ tends to $1$, but the interpolation constants explode; in practice libraries such as GMP switch from the schoolbook method to Karatsuba at a few dozen machine words, to Toom-$3$ and higher
variants beyond, and to FFT methods for thousands of words. The notebook measures these crossovers in Python.

\section{Fourier transforms over rings}\label{sec:ringfft}

For exact integer products, Tome~4, Theorem~7.11, computes the FFT modulo a prime $p\equiv1\pmod n$, which must be large enough to hold the coefficients of the product. Sch\"onhage and Strassen used a different ring, in which the roots of unity are powers of $2$, so that multiplying by them costs no multiplication at all.

\begin{definition}
Let $R$ be a commutative ring and $n\ge1$. An element $\omega\in R$ is a \emph{principal $n$th root of unity} if $\omega^n=1$ and $\sum_{j=0}^{n-1}\omega^{jk}=0$ for $0<k<n$.
\end{definition}

In a field, any element of order exactly $n$ is principal (the sum is $(\omega^{kn}-1)/(\omega^k-1)=0$); in a general ring, where $\omega^k-1$ may be a zero divisor, the condition must be checked.

\begin{proposition}\label{prop:ringdft}
Let $\omega$ be a principal $n$th root of unity in $R$, with $n$ invertible in $R$, and define the DFT by $\hat a_j=\sum_ka_k\omega^{-jk}$ for $a\in R^n$. Then $a_k=n^{-1}\sum_j\hat a_j\omega^{jk}$, and $\widehat{a*b}=\hat a\cdot\hat b$ for the cyclic convolution $(a*b)_k=\sum_la_lb_{k-l}$ (indices modulo $n$). If $n=2^r$ and
$\omega^{n/2}=-1$, the FFT of Tome~4, Theorem~7.9, computes the DFT with $\frac n2\log_2n$ multiplications by powers of $\omega$ and $n\log_2n$ additions.
\end{proposition}
\begin{proof}
$\sum_j\hat a_j\omega^{jk}=\sum_la_l\sum_j\omega^{j(k-l)}$, and the inner sum is $n$ if $l=k$ and $0$ otherwise, since $\omega^{-1}$ is principal too (the sum for $-k$ is the sum for $n-k$). The convolution identity is the computation of Tome~4, Theorem~7.6, which uses only the ring axioms and $\omega^n=1$. The FFT rests on Tome~4, Lemma~7.8, whose proof uses only
$\omega^{n/2}=-1$ and that $\omega^2$ is a principal $(n/2)$th root of unity; the latter holds because $\sum_{j<n/2}\omega^{2jk}=\frac12\sum_{j<n}\omega^{2jk}$ (each term appears twice), which vanishes for $0<k<n/2$, and $2$ is invertible since $n$ is.
\end{proof}

\begin{lemma}[Sch\"onhage and Strassen, 1971]\label{lem:fermatring}
Let $K$ be a power of $2$ and $R_K=\mathbb Z/(2^K+1)\mathbb Z$. Then $\omega=2$ is a principal $(2K)$th root of unity in $R_K$, $\omega^K=-1$, and $2K$ is invertible in $R_K$. Multiplying an element of $R_K$, stored as an integer in $[0,2^K]$, by a power of $2$ costs $O(K)$ bit operations.
\end{lemma}
\begin{proof}
$2^K\equiv-1$, so $2^{2K}\equiv1$. Let $0<k<2K$ and write $k=2^su$ with $u$ odd; then $2^s\le K$ since $2^s\le k<2K$ and both are powers of $2$. Let $d=K/2^s$; then $\omega^{dk}=2^{uK}\equiv(-1)^u=-1$, so the terms of $\sum_{j<2K}\omega^{jk}$ for $j$ and $j+d$ are opposite. Since $2d$ divides $2K$, the indices $j<2K$ split into blocks
$\{j,j+d\}$ with $j$ running over the $j$ with $\lfloor j/d\rfloor$ even, and the sum vanishes. The modulus is odd, so $2$, hence $2K$, is invertible. Finally, $2^s\cdot x$ with $s<2K$: write $2^sx=q2^K+r$ with $0\le r<2^K$, a shift; then $2^sx\equiv r-q$, and $q<2^{2K+1}$ splits in the same way, so $O(K)$ bit operations of shifting and subtracting suffice.
\end{proof}

\begin{remark}[the Sch\"onhage--Strassen algorithm, in outline]
To multiply two $N$-bit integers, cut them into $2^m\approx\sqrt N$ pieces of about $\sqrt N$ bits, viewed as polynomials evaluated at $2^{N/2^m}$. Their product polynomial has coefficients below $2^m2^{2N/2^m}$, so it can be computed in $R_K$ for $K$ a power of $2$ a little above $2N/2^m+m$, about $2\sqrt N$. In $R_K$, Lemma~\ref{lem:fermatring} gives a root of unity of
order $2K\ge2^{m+1}$, so the transforms cost $O(2^mm)$ operations each of $O(K)$ bits, that is $O(N\log N)$ bit operations, and the $2^m$ pointwise products are products in $R_K$, done recursively. (A twist by a root of unity of order $2^{m+1}$ turns the cyclic convolution into a negacyclic one, which is what reduction modulo $X^{2^m}+1$ gives, and halves the sizes.) The cost satisfies roughly
$\mathsf T(N)\le2^m\mathsf T(2\sqrt N)+O(N\log N)$; dividing by $N$, the cost per bit at most doubles while $\log N$ halves, and the $O(\log\log N)$ levels each contribute $O(\log N)$ per bit, hence $O(N\log N\log\log N)$. Harvey and van der Hoeven removed the last factor with multidimensional transforms and a careful choice of rings; their paper is the reading of this chapter.
\end{remark}

\section{Newton's method for power series}\label{sec:newton}

Power series over a commutative ring $R$ are written $f=\sum_{k\ge0}f_kx^k$; ``$g\equiv h\bmod x^n$'' means that $g$ and $h$ have the same first $n$ coefficients. A power series $f$ has an inverse $1/f$ exactly when $f_0$ is invertible in $R$, and its coefficients can be found one by one, in $O(n^2)$ operations for the first $n$. Newton's
iteration $g\mapsto g(2-fg)$, which approximates $1/f$ for real numbers using only multiplications, does much better.

\begin{theorem}[Sieveking, 1972; Kung, 1974]\label{thm:newtoninverse}
If $f_0$ is invertible, the first $n$ coefficients of $1/f$ can be computed with $O(\mathsf M(n))$ operations.
\end{theorem}
\begin{proof}
Suppose $fg\equiv1\bmod x^k$, so $1-fg=x^kh$ for a power series $h$. Then $f\cdot g(2-fg)=1-(1-fg)^2=1-x^{2k}h^2\equiv1\bmod x^{2k}$. So starting from $g=f_0^{-1}$ (correct modulo $x$) and setting $g\leftarrow g(2-fg)\bmod x^{2k}$ doubles the number of correct coefficients at each step. The step from $k$ to $2k$ needs $f\bmod x^{2k}$ times $g$ (degree less than $k$), then
times $g$ again, truncated: $O(\mathsf M(2k))$ operations. Summing over $k=1,2,4,\dots$ up to $n$ gives $O(\mathsf M(n))$ by the assumption on $\mathsf M$.
\end{proof}

The correct digits double, exactly as in Newton's method for $\sqrt2$, because the error $1-fg$ is squared at each step: here ``small'' means ``divisible by a high power of $x$''. The same principle computes square roots, logarithms and exponentials of power series in $O(\mathsf M(n))$ (Exercises~\ref{ex:sqrtnewton} and~\ref{ex:logseries}).

\begin{corollary}[fast division]\label{cor:fastdiv}
Let $a,b\in R[x]$ with $\deg a=n\ge\deg b=m$ and the leading coefficient of $b$ invertible. The quotient and remainder of $a$ by $b$ are computed with $O(\mathsf M(n))$ operations.
\end{corollary}
\begin{proof}
For a polynomial $p$ of degree at most $d$, let $\operatorname{rev}_d(p)=x^dp(1/x)$, its coefficients in reverse order. From $a=qb+r$ with $\deg q=n-m$ and $\deg r<m$, substituting $1/x$ and multiplying by $x^n$ gives $\operatorname{rev}_n(a)=\operatorname{rev}_{n-m}(q)\operatorname{rev}_m(b)+x^{n-m+1}\operatorname{rev}_{m-1}(r)$. So
$\operatorname{rev}_{n-m}(q)\equiv\operatorname{rev}_n(a)\cdot\operatorname{rev}_m(b)^{-1}\bmod x^{n-m+1}$, where $\operatorname{rev}_m(b)$ has constant coefficient the (invertible) leading coefficient of $b$. Theorem~\ref{thm:newtoninverse} gives the inverse modulo $x^{n-m+1}$, one product gives $\operatorname{rev}_{n-m}(q)$, hence $q$, and $r=a-qb$ costs one more product.
\end{proof}

For integers, the same Newton iteration on the reciprocal $1/b$, carried out with increasing precision, divides two $n$-bit integers in time $O(\mathsf M_{\mathbb Z}(n))$, where $\mathsf M_{\mathbb Z}$ is the cost of integer multiplication (Stephen Cook's thesis, 1966). Division, square roots and base conversion all cost a constant number of multiplications.

\section{Many points at once}\label{sec:multipoint}

Evaluating a polynomial of degree less than $n$ at one point costs $O(n)$ operations by Horner's rule, so at $n$ points $n^2$. The FFT does better for the special points $\omega^j$; for arbitrary points $u_0,\dots,u_{n-1}$ in a commutative ring, the trick is to divide. Let $n=2^k$ and define the \emph{subproduct tree}: its leaves are the
polynomials $m_{0,i}=x-u_i$, and the node at level $\ell$ is $m_{\ell,i}=m_{\ell-1,2i}\,m_{\ell-1,2i+1}=\prod_{j=i2^\ell}^{(i+1)2^\ell-1}(x-u_j)$; the root is $m=m_{k,0}=\prod_j(x-u_j)$ (Figure~\ref{fig:subproduct}).

\begin{figure}[ht]
\centering
\fig{ch09-subproduct}
\caption{The subproduct tree for $8$ points is built bottom-up by multiplications. The remainder tree runs through the same nodes top-down by divisions: each node receives $f$ modulo its product, and the leaves receive the values $f(u_i)=f\bmod(x-u_i)$.}
\label{fig:subproduct}
\end{figure}

\begin{theorem}[Fiduccia, 1972; Borodin and Moenck, 1974]\label{thm:multipoint}
The subproduct tree costs $O(\mathsf M(n)\log n)$ operations. Given it, a polynomial $f$ of degree less than $n$ can be evaluated at $u_0,\dots,u_{n-1}$ with $O(\mathsf M(n)\log n)$ operations.
\end{theorem}
\begin{proof}
At level $\ell$ there are $2^{k-\ell}$ products of two monic polynomials of degree $2^{\ell-1}$, costing $2^{k-\ell}\mathsf M(2^{\ell-1}+1)=O(\mathsf M(n))$ since $\mathsf M(n)/n$ is nondecreasing; there are $k=\log_2n$ levels. For evaluation, compute $r_{k,0}=f\bmod m_{k,0}=f$ and then, going down, $r_{\ell-1,2i+e}=r_{\ell,i}\bmod m_{\ell-1,2i+e}$ for $e=0,1$. By induction
$r_{\ell,i}=f\bmod m_{\ell,i}$: if $f=qm_{\ell,i}+r_{\ell,i}$, then since $m_{\ell-1,2i+e}$ divides $m_{\ell,i}$, $f$ and $r_{\ell,i}$ have the same remainder modulo $m_{\ell-1,2i+e}$. At the leaves, $f\bmod(x-u_i)=f(u_i)$. Each division at level $\ell$ divides a polynomial of degree less than $2^\ell$ by a monic one of degree $2^{\ell-1}$, which costs $O(\mathsf M(2^\ell))$ by Corollary~\ref{cor:fastdiv}; there are $2^{k-\ell+1}$ of them, so $O(\mathsf M(n))$ per level.
\end{proof}

\begin{theorem}[fast interpolation]\label{thm:interpolation}
Over a field, given distinct $u_0,\dots,u_{n-1}$ and values $v_0,\dots,v_{n-1}$, the unique polynomial $f$ of degree less than $n$ with $f(u_i)=v_i$ is computed with $O(\mathsf M(n)\log n)$ operations.
\end{theorem}
\begin{proof}
By Lagrange's formula, $f=\sum_iv_i\prod_{j\ne i}\frac{x-u_j}{u_i-u_j}=\sum_ic_i\frac m{x-u_i}$ with $c_i=v_i/\prod_{j\ne i}(u_i-u_j)$. The derivative of $m=\prod_j(x-u_j)$ is $m'=\sum_i\prod_{j\ne i}(x-u_j)$, so $m'(u_i)=\prod_{j\ne i}(u_i-u_j)$: the $n$ denominators are the values of $m'$ at the $u_i$, which Theorem~\ref{thm:multipoint} computes, and $n$ inversions give the $c_i$.
It remains to compute $\sum_ic_im/(x-u_i)$. For the node $(\ell,i)$ let $s_{\ell,i}=\sum_jc_jm_{\ell,i}/(x-u_j)$, the sum over the leaves $j$ below it. Then $s_{0,i}=c_i$, and splitting the leaves between the two children,
\[
s_{\ell,i}=s_{\ell-1,2i}\,m_{\ell-1,2i+1}+s_{\ell-1,2i+1}\,m_{\ell-1,2i},
\]
two products per node, $O(\mathsf M(n))$ per level going up. At the root, $s_{k,0}=f$.
\end{proof}

By the Chinese remainder theorem for polynomials (Tome~2, Chapter~6), evaluation at $n$ points is the isomorphism $K[x]/(m)\to K^n$ and interpolation its inverse; the same trees compute the Chinese remainder map for any coprime moduli, and for integers they give fast radix conversion and fast Chinese remaindering (Chapter~10).

\begin{think}
The subproduct tree costs $O(\mathsf M(n)\log n)$, but each evaluation at a single point costs only $O(n)$. Is the $\log n$ necessary for $n$ points? (Nobody knows: it is open whether multipoint evaluation can be done in $O(\mathsf M(n))$.)
\end{think}

\section{Matrices}\label{sec:matrices}

Tome~S1, Chapter~10, proved Strassen's identities and deduced that $n\times n$ matrices over any ring can be multiplied with $O(n^{\log_27})$ operations (Theorems~10.12 and~10.13), and defined the exponent $\omega$ of matrix multiplication, now known to be below $2.372$. Strassen's paper had a second title result: \emph{Gaussian elimination is not optimal}, because inversion reduces to multiplication.

\begin{theorem}[Strassen, 1969; Bunch and Hopcroft, 1974]\label{thm:inversion}
Suppose $n\times n$ real matrices can be multiplied with $\mathsf{MM}(n)=O(n^\tau)$ operations, for some $\tau>2$. Then an invertible $n\times n$ real matrix can be inverted, and its determinant computed up to sign, with $O(n^\tau)$ operations.
\end{theorem}
\begin{proof}
First let $A$ be symmetric positive definite, and $n=2^k$ (otherwise replace $A$ by $\begin{psmallmatrix}A&0\\0&I\end{psmallmatrix}$, still positive definite, at most doubling $n$). Write
$A=\begin{psmallmatrix}B&C^{\mathsf T}\\C&D\end{psmallmatrix}$ with blocks of size $n/2$. The block $B$ is positive definite (restrict the quadratic form to the first coordinates), hence invertible. Let $S=D-CB^{-1}C^{\mathsf T}$, the \emph{Schur complement}. For $y\ne0$, the vector $x=(-B^{-1}C^{\mathsf T}y,y)$ gives $x^{\mathsf T}Ax=y^{\mathsf T}Sy$, so $S$ is positive definite as well. Multiplying out shows
\[
A^{-1}=\begin{pmatrix}B^{-1}+B^{-1}C^{\mathsf T}S^{-1}CB^{-1}&-B^{-1}C^{\mathsf T}S^{-1}\\-S^{-1}CB^{-1}&S^{-1}\end{pmatrix},\qquad\det A=\det B\cdot\det S.
\]
So inverting $A$ takes two inversions of positive definite matrices of size $n/2$ ($B$ and $S$) and a constant number of products and sums of size $n/2$: $\mathsf I(n)\le2\mathsf I(n/2)+c\,\mathsf{MM}(n/2)+cn^2$. Unrolling, $\mathsf I(n)\le c\sum_{j\ge1}2^{j-1}\bigl(\mathsf{MM}(n/2^j)+(n/2^j)^2\bigr)+O(n)=O(n^\tau)$, the terms decreasing geometrically because $\tau>2$. The determinants follow along the same recursion.

For a general invertible $A$, the matrix $A^{\mathsf T}A$ is symmetric positive definite, and $A^{-1}=(A^{\mathsf T}A)^{-1}A^{\mathsf T}$: two more products. Also $\det(A^{\mathsf T}A)=(\det A)^2$ gives $|\det A|$.
\end{proof}

Over other fields, $A^{\mathsf T}A$ may be singular (over $\mathbb C$ use $A^*A$; over $\mathbb F_p$ it fails), and Bunch and Hopcroft's version of the recursion pivots, computing an LUP decomposition in $O(n^\tau)$; so inversion, determinants, and solving linear systems all cost $O(n^\omega)$ up to arbitrarily small exponent losses. Conversely, multiplication is no harder than inversion (Exercise~\ref{ex:mulviainv}), so $\omega$ is \emph{the} exponent of linear algebra.

\section{Powers and recurrences}\label{sec:recurrences}

To compute $x^e$ in a monoid, write $e$ in binary and square repeatedly: $x^{2j}=(x^j)^2$ and $x^{2j+1}=x\cdot(x^j)^2$. This uses at most $2\lfloor\log_2e\rfloor$ multiplications, an idea already in Pingala's Sanskrit treatise on prosody around 200 BCE, by most accounts, and in al-Kashi's work in 1427.

\begin{theorem}[Fiduccia, 1985]\label{thm:fiduccia}
Let $(u_k)$ satisfy $u_{k+d}=c_1u_{k+d-1}+\dots+c_du_k$ for all $k\ge0$, over a commutative ring, and let $P(x)=x^d-c_1x^{d-1}-\dots-c_d$. If $x^N\bmod P=\sum_{j<d}r_jx^j$, then $u_N=\sum_{j<d}r_ju_j$. Hence $u_N$ can be computed from $u_0,\dots,u_{d-1}$ with $O(\mathsf M(d)\log N)$ operations.
\end{theorem}
\begin{proof}
Let $L$ be the linear map on polynomials with $L(x^k)=u_k$. Then $L(x^kP)=u_{k+d}-c_1u_{k+d-1}-\dots-c_du_k=0$ for all $k$, so $L$ vanishes on all multiples of $P$. Writing $x^N=qP+\sum_jr_jx^j$, $u_N=L(x^N)=\sum_jr_ju_j$. To compute $x^N\bmod P$, square and multiply in $R[x]/(P)$: each step is a product of polynomials of degree less than $d$ followed by a division by the monic $P$, $O(\mathsf M(d))$ by
Corollary~\ref{cor:fastdiv}, and there are $O(\log N)$ steps.
\end{proof}

The textbook method raises the $d\times d$ companion matrix to the power $N$, which costs $O(d^\omega\log N)$; Fiduccia's polynomial version replaces $d^\omega$ by $\mathsf M(d)=O(d\log d)$. For the Fibonacci numbers, $P=x^2-x-1$ and the computation reduces to the doubling formulas of Exercise~\ref{ex:fibdouble}; Bostan and Mori showed in 2021 that a variant of the same idea, applied to the generating function, is faster still by a constant factor.

\begin{beauty}
One operation carries a whole field. Multiplication of polynomials, made fast by the roots of unity, becomes fast division by Newton's method, fast evaluation and interpolation by trees of products, fast inversion of matrices by the Schur complement, and fast linear recurrences by repeated squaring. The tower is built of
reductions, each a few lines of algebra, and every improvement to its foundation lifts the whole building.
\end{beauty}

\section{Exercises}

\begin{exercise}[$\star$]\label{ex:karatsubaint}
Show that Karatsuba's method multiplies two $n$-digit integers in $O(n^{\log_23})$ digit operations, although $a_0+a_1$ may have one more digit than $a_0$ and $a_1$.
\end{exercise}

\begin{exercise}[$\star\star$]\label{ex:toom3}
Write Toom-$3$ explicitly: with $c(Y)=C_0+C_1Y+C_2Y^2+C_3Y^3+C_4Y^4$ and the values $c(0),c(1),c(-1),c(2)$ and $C_4$, express $C_1,C_2,C_3$. Which divisions occur, and why does Toom-$3$ work over $\mathbb Z$ though $\mathbb Z$ is not a field?
\end{exercise}

\begin{exercise}[$\star\star$]\label{ex:sqrtnewton}
Over a field of characteristic not $2$, let $f$ be a power series with $f_0=1$. Show that the iteration $g\leftarrow\frac12(g+f/g)$, started at $g=1$, doubles the number of correct coefficients of $\sqrt f$, and that the first $n$ coefficients cost $O(\mathsf M(n))$.
\end{exercise}

\begin{exercise}[$\star$]\label{ex:logseries}
Over a field of characteristic $0$, let $f$ be a power series with $f_0=1$. Show that the first $n$ coefficients of $\log f$, defined as the power series with derivative $f'/f$ and constant term $0$, cost $O(\mathsf M(n))$.
\end{exercise}

\begin{exercise}[$\star$]\label{ex:mulviainv}
Show that $\begin{psmallmatrix}I&A&0\\0&I&B\\0&0&I\end{psmallmatrix}^{-1}=\begin{psmallmatrix}I&-A&AB\\0&I&-B\\0&0&I\end{psmallmatrix}$. Conclude that if $3n\times3n$ matrices can be inverted with $O(n^\tau)$ operations, $n\times n$ matrices can be multiplied with $O(n^\tau)$.
\end{exercise}

\begin{exercise}[$\star$]\label{ex:sqrt2}
In $R_K=\mathbb Z/(2^K+1)\mathbb Z$ with $4\mid K$, show that $2$ has order exactly $2K$ and that $\theta=2^{3K/4}-2^{K/4}$ satisfies $\theta^2=2$. What is the order of $\theta$, and why is it useful?
\end{exercise}

\begin{exercise}[$\star$]\label{ex:fibdouble}
Prove $F_{2k}=F_k(2F_{k+1}-F_k)$ and $F_{2k+1}=F_k^2+F_{k+1}^2$, and deduce an algorithm computing $F_N$ with $O(\log N)$ multiplications of integers. Relate it to Theorem~\ref{thm:fiduccia}.
\end{exercise}

\begin{exercise}[$\star\star$]\label{ex:taylorshift}
Over a field of characteristic $0$, show that the coefficients of $f(x+c)$, for $f$ of degree less than $n$, can be computed with $O(\mathsf M(n))$ operations. (Write $f(x+c)=\sum_kf_k(x+c)^k$ and expand.)
\end{exercise}

\begin{exercise}[$\star$]\label{ex:detabs}
Explain why Theorem~\ref{thm:inversion} only gives $|\det A|$ for a general real matrix, and how to recover the sign from the recursion applied to $A^{\mathsf T}A$ and one more computation.
\end{exercise}

\begin{exercise}[$\star\star$]\label{ex:remtree}
Show that evaluating a polynomial of degree less than $n$ at $n$ points by the remainder tree, with schoolbook multiplication and division, costs $O(n^2)$: the tree gains nothing without fast multiplication. With Karatsuba, what is the cost?
\end{exercise}

\section{Hints}

\begin{hint}{ex:karatsubaint}
$\mathsf T(n)\le3\mathsf T(\lceil n/2\rceil+1)+cn$. Compare with $\mathsf T'(m)=\mathsf T(m+2)$.
\end{hint}
\begin{hint}{ex:toom3}
$c(1)+c(-1)$ and $c(1)-c(-1)$ separate even and odd coefficients.
\end{hint}
\begin{hint}{ex:sqrtnewton}
If $g^2\equiv f\bmod x^k$, write $g^2-f=x^kh$ and compute $\bigl(\frac12(g+f/g)\bigr)^2-f$.
\end{hint}
\begin{hint}{ex:logseries}
Derivatives and primitives of power series cost $O(n)$.
\end{hint}
\begin{hint}{ex:mulviainv}
Multiply the two matrices.
\end{hint}
\begin{hint}{ex:sqrt2}
Use $2^K\equiv-1$ and expand $\theta^2$.
\end{hint}
\begin{hint}{ex:fibdouble}
Use $\begin{psmallmatrix}1&1\\1&0\end{psmallmatrix}^k=\begin{psmallmatrix}F_{k+1}&F_k\\F_k&F_{k-1}\end{psmallmatrix}$.
\end{hint}
\begin{hint}{ex:taylorshift}
The coefficient of $x^j$ in $f(x+c)$ is $\frac1{j!}\sum_kk!f_k\frac{c^{k-j}}{(k-j)!}$: a convolution.
\end{hint}
\begin{hint}{ex:detabs}
The sign of $\det A$ is not determined by $A^{\mathsf T}A$. Think of $A^{-1}$.
\end{hint}
\begin{hint}{ex:remtree}
Level $\ell$ has $2^{k-\ell}$ divisions of degree $2^{\ell}$ by degree $2^{\ell-1}$.
\end{hint}

\section{Solutions}

\begin{solution}{ex:karatsubaint}
\idea One extra digit per level changes the sizes by an additive constant, which the recursion absorbs.
\writeup Splitting $a=a_0+a_1\beta^{\lceil n/2\rceil}$ in base $\beta$, the three products have factors of at most $\lceil n/2\rceil+1$ digits, and the additions, subtractions and the final recombination cost $O(n)$ digit operations, so $\mathsf T(n)\le3\mathsf T(\lceil n/2\rceil+1)+cn$ for $n\ge4$. Let $\mathsf T'(m)=\mathsf T(m+3)$. For $m\ge1$,
$\lceil(m+3)/2\rceil+1\le\lceil m/2\rceil+3$, so $\mathsf T'(m)\le3\mathsf T'(\lceil m/2\rceil)+c(m+3)$ (with $\mathsf T$ nondecreasing, which we may assume by taking the maximum over smaller sizes), and the master theorem gives $\mathsf T'(m)=O(m^{\log_23})$, hence $\mathsf T(n)=O(n^{\log_23})$.
\end{solution}

\begin{solution}{ex:toom3}
\idea Evaluate at $0,\pm1,2,\infty$ and undo the Vandermonde system with small denominators.
\writeup $C_0=c(0)$ and $C_4$ are known. Let $s=\frac12(c(1)+c(-1))=C_0+C_2+C_4$ and $t=\frac12(c(1)-c(-1))=C_1+C_3$. Then $C_2=s-C_0-C_4$. Next, $c(2)=C_0+2C_1+4C_2+8C_3+16C_4$ gives $2C_1+8C_3=c(2)-C_0-4C_2-16C_4$, and with $2C_1+2C_3=2t$, $C_3=\frac16\bigl(c(2)-C_0-4C_2-16C_4-2t\bigr)$ and $C_1=t-C_3$.
The only divisions are by $2$ and by $6$ (that is $2$ and $3$), and they are exact: the $C_i$ are the coefficients of a product of polynomials with coefficients in $\mathbb Z$, so the quotients are integers. So Toom-$3$ works over $\mathbb Z$ with exact divisions, and over any ring where $2$ and $3$ are invertible. It uses $5$ products of a third of the size: $n^{\log_35}$.
\end{solution}

\begin{solution}{ex:sqrtnewton}
\idea Newton's step squares the error, here measured by powers of $x$.
\writeup If $g_0=1$ and $g^2-f=x^kh$, then $g$ is invertible, and $\frac12(g+f/g)=g-\frac{x^kh}{2g}$, so
$\bigl(\frac12(g+f/g)\bigr)^2-f=x^kh-x^kh+\frac{x^{2k}h^2}{4g^2}\equiv0\bmod x^{2k}$. Starting from $g=1$ (correct modulo $x$, as $f_0=1$), each step doubles the precision. The step to precision $2k$ needs $1/g$ modulo $x^{2k}$ (Theorem~\ref{thm:newtoninverse}, $O(\mathsf M(2k))$) and one product, so the total is $O(\mathsf M(n))$ by the doubling argument.
(Faster versions update $1/g$ by its own Newton step instead of recomputing it.)
\end{solution}

\begin{solution}{ex:logseries}
\idea Logarithm is integration of a quotient.
\writeup $f'$ costs $O(n)$, $1/f\bmod x^n$ costs $O(\mathsf M(n))$, one product gives $f'/f\bmod x^{n-1}$, and integrating termwise ($\sum a_kx^k\mapsto\sum\frac{a_k}{k+1}x^{k+1}$, possible in characteristic $0$) costs $O(n)$. Total $O(\mathsf M(n))$. The exponential is then obtained by Newton's method on the equation $\log g=f$: $g\leftarrow g(1+f-\log g)$.
\end{solution}

\begin{solution}{ex:mulviainv}
\idea Inversion of a block unipotent matrix produces a product as a byproduct.
\writeup Multiplying the two matrices, the diagonal blocks are $I$, the $(1,2)$ block is $-A+A=0$, the $(2,3)$ block is $-B+B=0$, and the $(1,3)$ block is $AB-AB+0=0$; so the product is $I$. Hence inverting one $3n\times3n$ matrix, built in $O(n^2)$, reveals $AB$ in its corner, and multiplication costs at most $\mathsf I(3n)+O(n^2)=O(n^\tau)$.
\end{solution}

\begin{solution}{ex:sqrt2}
\idea Square roots of $2$ give roots of unity of twice the order, for free.
\writeup $2^K=-1\ne1$ and $2^{2K}=1$, and the order divides $2K$ (a power of $2$) without dividing $K$, so it is $2K$. Then $\theta^2=2^{3K/2}-2\cdot2^K+2^{K/2}=2^K2^{K/2}-2^{K+1}+2^{K/2}\equiv-2^{K/2}+2+2^{K/2}=2$. So $\theta^{4K}=2^{2K}=1$ and $\theta^{2K}=2^K=-1$: $\theta$ has order $4K$. With it, transforms of length $4K$ in $R_K$ are possible, which allows splitting the
integers into twice as many pieces for the same ring; multiplying by $\theta$ costs two shifts and a subtraction.
\end{solution}

\begin{solution}{ex:fibdouble}
\idea Squaring the matrix power is doubling the index.
\writeup With $Q=\begin{psmallmatrix}1&1\\1&0\end{psmallmatrix}$, $Q^{2k}=(Q^k)^2$, and comparing entries: $F_{2k+1}=F_{k+1}^2+F_k^2$ and $F_{2k}=F_{k+1}F_k+F_kF_{k-1}=F_k(F_{k+1}+F_{k-1})=F_k(2F_{k+1}-F_k)$. So from $(F_k,F_{k+1})$ we get $(F_{2k},F_{2k+1})$ with three multiplications, and $(F_{2k+1},F_{2k+2})$ with one more addition; following the binary digits of $N$ from the top, $O(\log N)$ multiplications give $F_N$. In Theorem~\ref{thm:fiduccia}, $x^k\bmod(x^2-x-1)=F_kx+F_{k-1}$ (by induction), and squaring it modulo $x^2-x-1$ is exactly this doubling step.
\end{solution}

\begin{solution}{ex:taylorshift}
\idea A Taylor shift is a convolution in disguise, once factorials are pulled out.
\writeup $f(x+c)=\sum_kf_k\sum_j\binom kjc^{k-j}x^j$, so the coefficient of $x^j$ is $g_j=\frac1{j!}\sum_{k\ge j}(k!f_k)\frac{c^{k-j}}{(k-j)!}$. With $a_k=k!f_k$ and $b_i=c^i/i!$, $j!g_j=\sum_ka_kb_{k-j}$, which is a coefficient of the product of $\sum_ka_kx^{n-1-k}$ (reversed) and $\sum_ib_ix^i$. Computing $a$, $b$, one product and the divisions by $j!$ costs $O(\mathsf M(n))$.
\end{solution}

\begin{solution}{ex:detabs}
\idea $A^{\mathsf T}A$ forgets orientation; the sign needs $A$ itself.
\writeup $\det(A^{\mathsf T}A)=(\det A)^2$ is the same for $A$ and for $A$ with two columns swapped, whose determinant has the other sign; so no function of $A^{\mathsf T}A$ gives the sign. One way to get it within $O(n^\tau)$ is an LUP decomposition $A=PLU$ (Bunch and Hopcroft), where $\det A=\pm\det U$ with the sign of the permutation $P$. (In practice, LUP by Gaussian elimination with pivoting is the method.)
\end{solution}

\begin{solution}{ex:remtree}
\idea The tree only redistributes the work; the gain comes from fast products.
\writeup Building the tree with schoolbook products: level $\ell$ has $2^{k-\ell}$ products of degree $2^{\ell-1}$ polynomials, costing $O(4^{\ell-1})$ each, so $O(2^{k+\ell})$ per level, and the top levels dominate: $O(4^k)=O(n^2)$ in total. The divisions behave the same way. With Karatsuba, level $\ell$ costs $2^{k-\ell}O(3^\ell)=O(2^k(3/2)^\ell)$, again dominated by the top: $O(n^{\log_23})$, the cost of one multiplication of size $n$, without the $\log n$.
\end{solution}

\section*{Sources and further reading}

\begin{itemize}
\item \emph{A research paper to read.} D.~Harvey and J.~van der Hoeven, \emph{Integer multiplication in time $O(n\log n)$}, Annals of Mathematics 193 (2021), 563--617 (preprint posted on 18 March 2019). Read the introduction, which surveys the history from Karatsuba to F\"urer and states the conjecture of Sch\"onhage and Strassen; then the overview section, which explains how a one-dimensional transform of length $n$ is turned into a
multidimensional one with dimension about $\log n$ and why this avoids the $\log\log n$ cost of recursion. The details are demanding, but the architecture is exactly the reductions of this chapter.
\item \emph{The stories.} A.~Sch\"onhage and V.~Strassen, \emph{Schnelle Multiplikation gro\ss er Zahlen}, Computing 7 (1971), 281--292, with the conjecture in its last section; the announcement of 2019 is described by the University of New South Wales news release of April 2019 and by K.~Hartnett, \emph{Mathematicians discover the perfect way to multiply}, Quanta Magazine, April 2019.
\item A.~L.~Toom, \emph{The complexity of a scheme of functional elements realizing the multiplication of integers}, Soviet Mathematics Doklady 3 (1963), 714--716; S.~A.~Cook, \emph{On the minimum computation time of functions}, PhD thesis, Harvard University, 1966; D.~G.~Cantor and E.~Kaltofen, \emph{On fast multiplication of polynomials over arbitrary algebras}, Acta Informatica 28 (1991), 693--701; M.~F\"urer, \emph{Faster integer multiplication}, SIAM Journal on Computing 39 (2009), 979--1005.
\item M.~Sieveking, \emph{An algorithm for division of powerseries}, Computing 10 (1972), 153--156; H.~T.~Kung, \emph{On computing reciprocals of power series}, Numerische Mathematik 23 (1974), 341--348; C.~M.~Fiduccia, \emph{Polynomial evaluation via the division algorithm: the fast Fourier transform revisited}, Proceedings of the 4th ACM STOC (1972), 88--93;
A.~Borodin and R.~Moenck, \emph{Fast modular transforms}, Journal of Computer and System Sciences 8 (1974), 366--386; C.~M.~Fiduccia, \emph{An efficient formula for linear recurrences}, SIAM Journal on Computing 14 (1985), 106--112; A.~Bostan and R.~Mori, \emph{A simple and fast algorithm for computing the $N$-th term of a linearly recurrent sequence}, Symposium on Simplicity in Algorithms (SOSA 2021), 118--132.
\item V.~Strassen, \emph{Gaussian elimination is not optimal}, Numerische Mathematik 13 (1969), 354--356; J.~R.~Bunch and J.~E.~Hopcroft, \emph{Triangular factorization and inversion by fast matrix multiplication}, Mathematics of Computation 28 (1974), 231--236.
\item J.~von~zur~Gathen and J.~Gerhard, \emph{Modern Computer Algebra}, 3rd edition, Cambridge University Press, 2013, Chapters~8--10, is the reference for this chapter; D.~E.~Knuth, \emph{The Art of Computer Programming}, Volume~2, Section~4.3.3 and 4.6.3, for multiplication and powers, including the history of the binary method.
\end{itemize}
